\section{Pilot Methodology and Phases}
\label{sec:pilot-methodology}

\subsection{Overview}

The pilot follows a \textbf{gated, phased} methodology aligned with safety science and
human-subjects research norms.
No phase involving live occupant sensing begins without documented governance approval.

\subsection{Phase 0: Requirements and Institutional Alignment (Current)}

Activities:
\begin{itemize}[noitemsep]
  \item distribute this requirements document to stakeholders;
  \item site walkthrough and zone mapping;
  \item IT capacity assessment;
  \item initiate DPIA and ethics review;
  \item protocol mapping workshop (planned).
\end{itemize}

Exit criteria: signed intent to proceed; steering committee formed.

\subsection{Phase 1: Sandbox and Tabletop (No Live Public Sensing)}

Activities:
\begin{itemize}[noitemsep]
  \item deploy infrastructure on isolated VLAN with synthetic feeds;
  \item train safety officers on dashboard;
  \item tabletop exercises with scripted semantic events;
  \item refine protocol mapping matrix.
\end{itemize}

Duration: 4--8 weeks (indicative).

Exit criteria: training complete; tabletop debrief documented; DPIA draft reviewed.

\subsection{Phase 2: Minimum Viable Pilot (MVP)}

\textbf{Scale:} 1 school, 2--4 zones, 1--2 edge nodes, single orchestrator.

Activities:
\begin{itemize}[noitemsep]
  \item limited live sensing in approved zones only;
  \item staffed-hours monitoring with defined review SLAs;
  \item collect evaluation metrics (Section~\ref{sec:evaluation-metrics});
  \item weekly calibration and false-alarm review.
\end{itemize}

Duration: 8--16 weeks (indicative).

Exit criteria: governance review; decision to expand, modify, or halt.

\subsection{Phase 3: Recommended Pilot}

\textbf{Scale:} 1--3 schools, 6--12 zones total, heterogeneous zone types (hallway, cafeteria, exit).

Activities:
\begin{itemize}[noitemsep]
  \item cross-site comparison of operator workflows;
  \item optional GPU video nodes at one site;
  \item structured staff feedback (surveys/interviews with ethics approval);
  \item prepare Horizon Europe or national funding evidence pack.
\end{itemize}

Duration: 6--12 months (indicative).

\subsection{Phase 4: Extended Pilot}

\textbf{Scale:} 5--10 schools.

Activities:
\begin{itemize}[noitemsep]
  \item multi-site federation patterns;
  \item hardened bus security (mTLS);
  \item optional anonymized metrics export to research server;
  \item pre-registered inferential evaluation if ethics-approved.
\end{itemize}

Duration: 12--24 months (indicative, funding dependent).

\subsection{Minimum Viable Pilot Configuration}

\begin{table}[htbp]
  \centering
  \caption{MVP minimum configuration.}
  \label{tab:mvp-config}
  \small
  \begin{tabular}{@{}ll@{}}
    \toprule
    Element & MVP minimum \\
    \midrule
    Schools & 1 \\
    Zones & 2--4 \\
    Edge nodes & 1--2 (Edge-S or Edge-V) \\
    Event bus & 1 NATS/MQTT broker \\
    Orchestrator & 1 host (32 GB RAM) \\
    Dashboard & 1 API + web UI \\
    GPU & Optional (CPU-only acceptable for MVP) \\
    Research server & Disabled \\
    Staff FTE (school) & 0.2--0.5 safety officer; 0.1 IT \\
    \bottomrule
  \end{tabular}
\end{table}

\subsection{Go / No-Go Gates}

\begin{enumerate}[noitemsep]
  \item Governance charter signed.
  \item DPIA residual risk accepted by DPO and leadership.
  \item Protocol mapping matrix approved.
  \item Training completion $\geq$ 90\% of designated staff.
  \item Tabletop exercises completed with documented remediations.
  \item IT sign-off on VLAN and firewall rules.
  \item Ethics approval for evaluation instruments (if applicable).
\end{enumerate}

\subsection{Stop Conditions}

Immediate pause triggers:
\begin{itemize}[noitemsep]
  \item privacy policy violation or raw media egress detected;
  \item unauthorized access to dashboard or audit logs;
  \item stakeholder withdrawal or regulatory inquiry requiring halt;
  \item unmitigated false-alarm rate exceeding pre-agreed pilot threshold.
\end{itemize}

Stop conditions invoke incident review before resumption.
